\section{Conclusion}

NephroVision is an end-to-end biomedical AI prototype for automated kidney and renal tumor/cyst segmentation from contrast-enhanced abdominal CT volumes. The system integrates CT preprocessing (HU clipping $[-200, 300]$, min--max normalization), a 3D U-Net for volumetric semantic segmentation, sliding-window inference ($64 \times 192 \times 192$ patches, $50\%$ overlap), 8-flip test-time augmentation, conservative connected-component post-processing, volumetric statistics computation, and interactive browser-based 3D visualization. Comprehensive development documentation records preprocessing decisions, training and inference configurations, engineering choices, challenges, failure cases, and reproducibility information. Safety documentation positions the system as IEC 62304 Software Safety Class B decision-support requiring physician review.

Evaluated on 64 independent held-out cases from the KiTS23 dataset, NephroVision achieved a kidney Dice of $0.9307 \pm 0.064$, tumor Dice of $0.6558 \pm 0.262$, a mean Dice of $0.7933$, and a tumor detection rate of $100\%$ ($64/64$). Kidney segmentation is robust and consistent, with clear CT boundaries and stable morphology supporting reliable delineation across diverse cases. Tumor and cyst detection is reliable, with every tumor-containing case correctly identified. Tumor boundary accuracy remains moderate, reflecting the inherent difficulty of segmenting small, heterogeneous, and poorly-marginated lesions on single-phase contrast-enhanced CT.

The project demonstrates that an academic team can build a functional medical image segmentation system with documented engineering practices. The honest reporting of moderate tumor boundary precision---distinguishing detection from segmentation---strengthens rather than weakens the project's credibility. The decision-support framing, explicit limitations documentation, and intended-use boundaries ensure that the system is positioned appropriately as an academic prototype, not as a clinically validated device.

Future work should prioritize improving tumor boundary accuracy through advanced loss functions or stronger frameworks (nnU-Net), external multi-center validation to assess domain shift and generalizability, and prospective clinical evaluation with radiologist reviewers. Longer-term preparation for clinical and regulatory readiness---including formal risk management, cybersecurity assessment, usability evaluation, and quality management system development---would be required before any clinical deployment.

NephroVision is not clinically approved. All outputs require physician review. The system does not replace radiologists---it provides segmentation support that can assist volumetric assessment while keeping the physician as the decision-maker.
